\FloatBarrier
\section{Deployment, economics and limitations}\label{sec:deployment}

\subsection{Operational numerical weather prediction}
An operational centre runs a cycle every six hours: an analysis, a forecast ensemble, and a
product generation step. \ac{1} is sized for the forecast and competent at the analysis. A
single card holds a 64-member, 4.19 M-cell, 24-variable ensemble and completes a 24-hour forecast
in two minutes at 57 W, which means a four-card node can run the same cycle for a domain four
times larger, or run four independent domains, at 260 W. The equivalent H100 node draws 1.05 kW at
idle and 2.8 kW under load for the same throughput.

The scale-out story is the weakest part of the design. \ac{1} has one PCIe 6.0 x16 link and no
cache-coherent die-to-die fabric. A regional domain that fits in 48 GB is a single-card job; a
global ensemble is not, and the halo traffic between cards would go through PCIe at 100 GB/s,
which is 24$\times$ slower than the on-card fabric. A production cluster would therefore
decompose by ensemble member (each card owns a subset of members for the full domain) rather than
by spatial subdomain, which keeps inter-card traffic to boundary conditions and output. That is a
real constraint and it shaped the design: the machine is a single-domain appliance, not a
scale-out building block.

\subsection{Research groups and forecast-as-a-service}
The second deployment is a university group or a commercial provider that runs many small
ensembles rather than one large one. Here the relevant numbers are the efficiency mode and the
cost of ownership. At 1.05 GHz and 0.55 V the card delivers the same forecast in 172 s at 33 W,
which is a 20\% energy reduction for a 37\% time penalty; for a job that runs overnight, that is
free. The five-year cost per forecast day is \$0.41 against \$8.22 for the H100, and the gap is
capital cost, not electricity.

\subsection{Emergency response and dispersion}
A third deployment is time-critical and safety-relevant: a release of a hazardous substance, a
volcanic ash cloud, a wildfire smoke plume. The requirement is not throughput but a bounded
response time and a reproducible answer that can be defended after the event. \ac{1} has a
static schedule, a compiler-known cycle count per tile and no dynamic scheduling in the numerical
path, so the worst-case execution time of a dispersion calculation is a computed number rather
than a statistical one. The same property gives bit-reproducibility, which matters when a
regulator asks why the plume was predicted to go east rather than west.

\subsection{Limitations and threats to validity}
The report is only useful if its weaknesses are stated as clearly as its results.

\begin{enumerate}[leftmargin=1.4em,itemsep=1pt]
\item \textbf{No silicon.} Every chip number is a model projection. The model is calibrated
      against published data and its assumptions are sampled, but a first silicon part can
      plausibly land 20--30\% away from the projected energy on real kernels.
\item \textbf{The GPU baseline is a model.} The H100 comparison assumes a 0.30 stencil efficiency
      and a particular power state. Both are defensible and both are varied; the conclusion
      survives at 100\% GPU utilisation with a 3.9$\times$ ratio, but a reader who believes GPUs
      achieve 60\% on this workload should read the 8.2$\times$ end of Fig.~\ref{fig:cuda}(e).
\item \textbf{The workload is a subset.} Chemistry beyond passive tracers, full radiation with
      hundreds of bands, atmospheric composition and land-surface models are not in the
      benchmark. They are more SFU-heavy, which would help the ASIC's relative position, but they
      are not measured.
\item \textbf{Precision validation is on a surrogate.} Lorenz--96 supports the argument, not the
      design point. A full-model study is the first bring-up task.
\item \textbf{No scale-out study.} The inter-card fabric is a PCIe link and the report does not
      pretend otherwise.
\item \textbf{No software maturity claim.} The compiler described in Section~\ref{sec:software}
      is a design, and the application that ships with this report runs on a virtual device.
\end{enumerate}

\subsection{What would falsify the thesis}
A single measurement would settle most of the argument: run the ten kernels of Table~\ref{tab:kernels}
on an H100 with a state-of-the-art stencil compiler and measure joules per delivered forecast at
matched skill. If that number is within 2$\times$ of the projection in this report, the
architectural thesis is wrong and the correct conclusion is that CUDA is good enough. The
prediction of this report is that it will be 6--10$\times$ worse, and the prediction is stated in
those terms so that it can be checked.

\begin{table}[t]
\centering\small
\caption{Deployment comparison. Energy is per forecast-and-analysis day; cost is five-year total
cost of ownership per day at \$0.12/kWh.}
\label{tab:deploy}
\begin{tabular}{@{}lrrr@{}}
\toprule
Scenario & ACPU-1 node & H100 node & Ratio\\
\midrule
Cards & 4 & 4 & \\
Idle power & 24 W & 360 W & 15$\times$\\
Load power (W1) & 228 W & 1.05 kW & 4.6$\times$\\
Forecast time (same domain) & 126 s & 278 s & 2.2$\times$ faster\\
Energy per forecast day & 34.2 kJ & 316 kJ & 9.2$\times$\\
Capital cost & \$6.0k & \$120k & 20$\times$\\
Cost per forecast day & \$1.64 & \$32.9 & 20$\times$\\
\bottomrule
\end{tabular}
\end{table}
